\section{The integral Lefschetz defect}\label{sec:lattice}

We first work with a unimodular symplectic lattice \(\Lambda\) of rank
\(2g\), \(g\geq4\), and polarization \(\theta\).  This isolates the only
integral linear algebra in the paper.

\begin{theorem}\label{thm:symplectic-saturation}
For
\[
 L_3:\bigwedge^3\Lambda\longrightarrow\bigwedge^5\Lambda,
 \qquad \alpha\longmapsto\theta\wedge\alpha,
\]
one has
\begin{equation}\label{eq:general-saturation}
 \Sat(\im L_3)=im L_3+\theta^{[2]}\wedge\Lambda.
\end{equation}
Moreover
\begin{equation}\label{eq:tau-isomorphism}
 \tau:\Lambda/2\Lambda\xrightarrow{\sim}
 \Sat(\im L_3)/\im L_3,
 \qquad z\longmapsto[\theta^{[2]}\wedge z]
\end{equation}
is symplectically natural.  The nonzero Smith factors of \(L_3\) are all
one except for exactly \(2g\) factors equal to two.
\end{theorem}

\begin{proof}
Choose a symplectic basis \(e_1,f_1,\ldots,e_g,f_g\), and put
\(x_i=e_i\wedge f_i\).  Label an exterior monomial by the symplectic pairs
that occur as singletons and by those that occur as full pairs.  Fix a
singleton set \(A\), choose \(e_i\) or \(f_i\) for each \(i\in A\), and
write \(u_A\) for their product.  For \(T\cap A=\varnothing\), put
\(x_T=\prod_{i\in T}x_i\).  Multiplication by
\(\theta=\sum_i x_i\) preserves the singleton label and acts by
\begin{equation}\label{eq:up-incidence}
 u_Ax_T\longmapsto
 \sum_{j\notin A\cup T}u_Ax_{T\cup\{j\}}.
\end{equation}
These are integral direct-sum blocks; no rational primitive decomposition is
being used.

In degree three there are two cases.  If \(|A|=3\), then \(T\) is empty,
and \eqref{eq:up-incidence} sends one basis vector to an all-ones vector.
Its image is primitive.  If \(|A|=1\), then \(|T|=1\), and the block is the
unsigned vertex--edge incidence map of the complete graph on
\(n=g-1\) vertices:
\[
 U_n:\Z^n\longrightarrow\Z^{\binom n2},
 \qquad [i]\longmapsto\sum_{j\ne i}[\{i,j\}].
\]

For \(n\geq3\), \(U_n\) is injective: the equations
\(a_i+a_j=0\) on a triangle force every \(a_i\) to vanish.  An
\((n-1)\)-minor is one, using the rows \(\{1,j\}\), \(2\leq j\leq n\),
and columns \(2,\ldots,n\).  Every \(n\)-minor is even because modulo two
the sum of all columns is zero.  Finally, the rows of one triangle together
with \(\{1,k\}\), \(4\leq k\leq n\), give an \(n\)-minor of determinant
\(\pm2\).  Hence
\begin{equation}\label{eq:complete-graph-smith}
 \operatorname{SNF}(U_n)=\operatorname{diag}(1^{n-1},2).
\end{equation}
The order-two class is represented by the sum of all edge vectors: twice
that vector is the sum of all vertex columns, while the triangle equations
show that it is not itself in the image.

There is one such block for each of the \(2g\) choices of a singleton
basis vector.  Since
\[
 \theta^{[2]}=\sum_{i<j}x_i\wedge x_j,
\]
the vector \(\theta^{[2]}\wedge e_k\), or
\(\theta^{[2]}\wedge f_k\), is precisely the all-edge vector in the
corresponding block.  The blocks are independent, so these \(2g\) vectors
exhaust the saturation defect and prove \eqref{eq:general-saturation} and
\eqref{eq:tau-isomorphism}.

Although the calculation used a basis, \(\tau\) is defined by the
polarization and its integral divided square.  It therefore commutes with
every element of \(\operatorname{Sp}(\Lambda,\theta)\).
\end{proof}

For \(g=5\), the domain has rank \(\binom{10}{3}=120\).  Theorem
\ref{thm:symplectic-saturation} gives ten factors of two and 110 factors of
one, proving the algebraic part of \eqref{eq:main-saturation}.  Abstractly,
these Smith factors and the quotient \(\Lambda/2\Lambda\) are the
off-central specialization of the integral Lefschetz filtration and
compatible splittings in
\cite[Theorem~2.9, Corollary~2.10, and Proposition~2.14]{FaulknerMiller}.
The occupancy proof above gives a direct block calculation and identifies
the divided-power representatives \(\theta^{[2]}\wedge z\) that enter the
geometric endpoint.
